\chapter{功能测试}\label{ch:func}

本章按选题指南的 5 项功能报告测试结果，最后给出 5 项功能在同一次运行中的联合测试。各节先列测试目的，再以测试项表给出每个用例的配置、激励、判据和结果，并附关键时序的波形。

\section{ARINC 664 Part 7 通信}\label{sec:f-comm}

\textbf{测试目的：}验证以太网 MAC 收发、UDP/IP 封装、VL 寻址、发送缓存与 BAG 调度、双网发送、接收映射与端口分发、配置口和复位恢复，对应测试点 F1-01 至 F1-14。

\begin{items}{ARINC 664 Part 7 通信测试项}{tab:f-comm}
BAS-01 & 四种宽度和深度的存储器配置；随机写入与同步读，含最低、最高地址，禁写期间给出不同数据 & 读出数据与写入一致；同步读延迟正确；读使能为低时输出保持 & 20764/0\newline --- & 通过 \\
MAC-01 & 消息长度 1、17、18、1471 字节及随机长度；网络使能取仅 A、仅 B、双网、全关；首个请求覆盖 8 种时钟相位；在帧首、帧中、帧尾断供上游字节 & 前导码、SFD、填充、FCS 正确；帧间隔 12 拍；断供时输出错误 FCS，不截断帧 & 264649/0\newline --- & 通过 \\
MAC-02 & 正常帧，以及截短帧、超长帧、坏前导码、坏 FCS；每个异常帧后紧跟一个正常帧 & 帧边界正确，去 FCS 后输出；异常帧的错误分类正确；随后的正常帧正确接收 & 105303/0\newline --- & 通过 \\
MAC-03 & 发送输出接回接收；不同帧长；A、B 两网 & 每网接收内容与发送内容一致 & 1385254/0\newline --- & 通过 \\
\midrule
TX-01 & 默认缓存（256 页 × 512 字节）；消息长度 1、16、17、18、255--257、511--513、1025、1470、1471 字节及 32 个随机长度 & 线路帧与参考模型逐字节一致，含 VLID、网络标识、IP/UDP 头、端口、载荷、填充、SN、FCS；SN 复位后首帧为 0，255 后回绕到 1 & 662055/0\newline 7605/0 & 通过 \\
TX-02 & 缓存改为 32 页 × 2048 字节；激励同 TX-01 & 同 TX-01 & 661635/0\newline 7605/0 & 通过 \\
TX-03 & 改变 BAG、消息相对允许时刻的提交时间和 BAG 计数初值；覆盖持续排队、晚到消息、计数回绕 & 同一 VL 相邻两帧的起点间隔不小于 BAG；帧内容正确 & 14054170/0\newline 14591/0 & 通过 \\
TX-04 & 一帧长帧占用线路期间，向全部 128 条 VL 同时提交消息 & 各 VL 队列独立；按轮询顺序发出；网络选择与配置一致；帧内容正确 & 14438575/0\newline 56211/0 & 通过 \\
TX-05 & 构造 VL 禁用、队列满、消息超长、无空闲页四类拒绝，随后不复位重新填满缓存 & 拒绝类别与计数正确；被拒消息已占用的页全部归还 & 30267256/0\newline 4642/0 & 通过 \\
TX-06 & 固定输入向量 & 发送记录与参考一致；记录作为网表对照基准 & 219313/0\newline 229/0 & 通过 \\
\midrule
RX-01 & 改变消息长度、UDP 端口、帧间隔；VLID 取 0 与 65535，接收索引取 0 与 4095；配置全部 4096 条接收 VL，每条收两对双网副本 & 映射与端口分发正确；交付 8192 条、冗余丢弃 8192 个，逐条一致 & 9960369/0\newline 123855/0 & 通过 \\
\midrule
SYS-01 & 每端 4 条 VL，BAG、最大长度、网络选择各不相同；双向互发长短消息 & 每次交付对应对端一条被接纳的消息，内容、VL、端口、顺序正确；事件与交付一致 & 5155148/0\newline 2896/0 & 通过 \\
SYS-04 & 分两轮用满 128 条发送 VL；接收端配置 4096 条，抽查含 0、4095 的 256 个索引；构造发送侧四类拒绝 & 交付正确；同一 VL 发送间隔不小于 BAG；拒绝计数正确；各 VL 配置互不影响 & 216810445/0\newline 116677/0 & 通过 \\
SYS-05 & 配置口读写；业务进行中改写配置；改小 VL 最大长度 & 读结果在请求后第 2 拍返回；计数值为请求前一拍的累计数；写握手完成；改小后长消息被拒、短消息正常接纳 & 21167623/0\newline 714342/0 & 通过 \\
SYS-06 & 对端连续发送 8 条最长消息期间复位接收端，重新初始化和配置后继续发送 & 复位前的残帧不进入新的交付；复位后交付正确 & 10991314/0\newline 3682/0 & 通过 \\
SYS-09 & eLinx 界面入口；两端各发送 1、64、1471 字节消息 & 端口逐字节一致；双网副本只交付一次 & \ElinxChecks/0\newline --- & 通过 \\
\end{items}

\subsection{发送侧拒绝与页回收}

SYS-04 在端系统 1 上构造发送侧四类拒绝，经配置口读回计数（表~\ref{tab:fault-tx}）。

\begin{table}[!htbp]\centering
\caption{发送侧拒绝的注入与读回}\label{tab:fault-tx}
\begin{tabularx}{\linewidth}{@{}lLl@{}}
\toprule
拒绝类别 & 注入方式 & 计数器读回 \\
\midrule
VL 禁用 & 向禁用的 VL 125 提交 1 条消息 & \code{0x41} = 1 \\
队列满 & 向队列限额为 0 的 VL 127 提交 1 条消息 & \code{0x42} = 1 \\
消息超长 & 向最大长度为 1 的 VL 126 提交 18 字节消息 & \code{0x43} = 1 \\
无空闲页 & 85 条 1471 字节消息占用 255 页后，依次提交 1471、512、1 字节消息 & \code{0x44} = 2 \\
\bottomrule
\end{tabularx}
\end{table}

无空闲页一项同时检验页回收。剩余 1 页时，1471 字节消息需要 3 页，在申请第 2 页时失败，已取得的 1 页归还；随后的 512 字节消息使用该页，发送成功；最后的 1 字节消息因页池为空被拒绝。三条消息的结果依次为拒绝、发送、拒绝，与页回收后的预期一致。

\subsection{仿真波形}

图~\ref{fig:wtx} 为 ModelSim 原生窗口截取的应用消息双网发送仿真波形（\code{sim/wave/png/tx_dual_*.png}）。子图 (a) 显示应用口握手写入 1 字节消息及发送调度过程；子图 (b) 放大显示 A、B 两网 GMII 接口同拍发出前导码和帧数据的时序。

\figWaveDual[应用消息与双网同拍发送]{tx_dual_1}{tx_dual_2}{宏观时序：应用口写入与发送调度}{微观时序：A、B 两网 GMII 同拍发出前导码与数据}{ModelSim 仿真波形：应用消息与双网同拍发送}{fig:wtx}

\section{高完整性通信}\label{sec:f-ic}

\textbf{测试目的：}验证帧错误检测、完整性检查（IC）、冗余管理（RM）、接收缓存溢出处理和错误计数，对应测试点 F2-01 至 F2-08。

\begin{items}{高完整性通信测试项}{tab:f-ic}
RX-02 & 输入 FCS 错、帧长错、IP 头错、UDP 头错、未知 VL、网络不允许的帧 & 每帧拒绝原因与参考判决一致；无交付 & 576923/0\newline 890/0 & 通过 \\
RX-03 & 序号正常递增、丢一帧、重复、跳号、复位后 0 号帧、回绕；到达次序取 A 先、B 先、同拍 & 接纳与拒绝符合 IC 规则；先到副本交付，后到副本丢弃 & 4635288/0\newline 1274/0 & 通过 \\
RX-04 & 64 条 VL 混合收发长短消息；IC、RM 四种开关组合 & 各 VL 序号状态互不影响；去重正确 & 6406189/0\newline 19180/0 & 通过 \\
RX-05 & 字节环缩为 512 字节；长消息与短消息交替 & 存不下的帧被丢弃并计数，暂占空间释放；后续帧正常接收 & 422837/0\newline 99/0 & 通过 \\
RX-06 & 交付队列缩为 2 项；输入 8 帧 & 队满事件 1 次、冗余丢弃 2 次、交付 5 条；释放空间后正常交付 & 373852/0\newline 57/0 & 通过 \\
RX-07 & 固定的正常及异常输入向量 & 接收记录与参考一致；记录作为网表对照基准 & 374111/0\newline 110/0 & 通过 \\
RX-08 & 两网间隔 0--20 拍；帧尾位置在计时边界前后各 32 拍；交换领先网络；接收引擎忙时推进时间 & 冗余窗口内只交付一次；过期记录失效后同一序号可再次交付 & 353130/0\newline --- & 通过 \\
\midrule
SYS-02 & 链路模型在线路上改字节、插入错误帧、重放、延后重放（表~\ref{tab:fault-rx}） & 各类计数器读回与注入数一一对应；正常消息交付不受影响 & 47117674/0\newline 8726/0 & 通过 \\
SYS-03 & 两网延迟取相等、A 晚、B 晚、恰在冗余窗口边界 & 窗口内只交付一次，窗口外两份各自交付 & 15481734/0\newline 3763/0 & 通过 \\
SYS-07 & 字节环缩为 512 字节；每端发送 513、1000、1471 字节消息，与短消息交替 & 每端溢出计数 3/3；同时正常交付 3 条短消息 & 9681100/0\newline 1488/0 & 通过 \\
\end{items}

\subsection{异常注入与计数读回}

接收侧共注入 10 类异常，与发送侧 4 类拒绝（表~\ref{tab:fault-tx}）合计 14 类。每个阶段结束后等待通路排空，再读统计寄存器，与该阶段的注入记录核对；计数器在同一次运行中累计，部分行按本阶段增量核对。表~\ref{tab:fault-rx} 中计数器地址为十六进制，“A/B”指两网计数，“每端”指两个端系统各自的读回值。

\begin{longtable}{@{}>{\raggedright\arraybackslash}p{2.0cm}>{\raggedright\arraybackslash}p{1.3cm}>{\raggedright\arraybackslash}p{5.6cm}>{\raggedright\arraybackslash}p{5.7cm}@{}}
\caption{接收侧异常的注入与读回}\label{tab:fault-rx}\\
\longtablehead{4}{异常类别 & 用例 & 注入方式与数量 & 计数器读回}
FCS 错误 & SYS-02 & 改写字节，每端每网 2 帧，共 8 帧 & 每端 \code{0x49/0x51} 由 0 增至 2/2 \\
帧长错误 & SYS-02 & 插入错误帧，每端每网 1 帧，共 4 帧 & 每端 \code{0x4a/0x52} = 1/1 \\
IP 头错误 & SYS-02 & 插入错误帧，每端每网 1 帧，共 4 帧 & 每端 \code{0x4b/0x53} = 1/1 \\
UDP 头错误 & SYS-02 & 校验和字段非零，每端每网 1 帧，共 4 帧 & 每端 \code{0x4b/0x53} = 2/2（与 IP 头共用计数器，本阶段各加 1） \\
未知 VL & SYS-02 & 每端每网 1 帧，共 4 帧 & 每端 \code{0x4c/0x54} = 1/1 \\
网络不允许 & SYS-02 & 每端 2 条双网消息，A、B 各 1 个副本不被允许，共 4 个 & 每端 \code{0x4d/0x55} = 1/1 \\
同网重复 & SYS-02 & 四条单向链路各重放 1 帧，共 4 帧 & 每端 \code{0x4e/0x56} = 1/1 \\
序号乱序 & SYS-02 & 四条单向链路各延后重放 1 帧，共 4 次 & 每端 \code{0x4e/0x56} = 2/2（本阶段各加 1） \\
字节环溢出 & SYS-07 & 字节环 512 字节，每端发送 513、1000、1471 字节消息各 1 条 & 每端 \code{0x4f/0x57} = 3/3 \\
交付队列满 & RX-06 & 队列 2 项，输入 8 帧 & 队满事件 1 次、冗余丢弃 2 次、交付 5 条 \\
\end{longtable}

同网重放的帧序号重复，由 IC 拒绝并计入完整性错误；双网冗余副本通过 IC 后由 RM 丢弃，计入冗余丢弃。两种重复的处理路径和计数器相互独立。正常冗余场景中，SYS-04 的端系统 2 收到 922 组双网消息共 1844 个副本，冗余丢弃计数 \code{0x58} 与交付计数 \code{0x59} 均为 922。

\subsection{仿真波形}

图~\ref{fig:wrm} 为 ModelSim 原生窗口截取的双网副本去重仿真波形（\code{sim/wave/png/rx_redundancy_*.png}）：A 网副本先到并交付应用口（子图 a）；晚到 2002 拍的 B 网副本到达后被冗余管理（RM）模块丢弃，触发 \code{ev_rm_drop} 脉冲（子图 b）。

\figWaveDual[双网副本接收与去重]{rx_redundancy_1}{rx_redundancy_2}{A 网副本先到：完整性检查通过并向应用口交付}{晚到 2002 拍的 B 网副本到达：触发冗余丢弃 ev\_rm\_drop}{ModelSim 仿真波形：两网副本先后到达与冗余去重}{fig:wrm}

图~\ref{fig:wfault} 为 FCS 错误帧的拒绝与配置口读回波形（\code{sim/wave/png/rx_fcs_error_*.png}）：坏 FCS 帧到达后触发 \code{ev_a_fcs_error} 脉冲，应用口无交付（子图 a）；随后的独立配置读请求经两拍流水线返回 \code{count_a_fcs = 1}（子图 b）。

\figWaveDual[FCS 错误帧的拒绝与计数]{rx_fcs_error_1}{rx_fcs_error_3}{坏 FCS 帧到达：触发 ev\_a\_fcs\_error，应用口不交付}{配置口读回：经两拍返回 count\_a\_fcs = 1}{ModelSim 仿真波形：FCS 错误帧的拒绝与计数器读回}{fig:wfault}

\section{SNMP 网络管理}\label{sec:f-snmp}

\textbf{测试目的：}验证 SNMP 代理的查询、设置、拒绝和通知功能，对应测试点 F3-01 至 F3-06。管理端发出的请求报文经 AFDX 通路到达被测端代理；请求与响应均由第三方库 \code{pyasn1} 编解码，响应中的数值与配置口读回值比较。

\begin{items}{SNMP 测试项}{tab:f-snmp}
SVC-01 & Get、GetNext；v1、v2c；逐对象遍历至表尾；错误团体名；请求号正负值；BER 长度边界（16 字节团体名、长形式长度）；穿插双向用户消息 & 响应字段经独立解码正确；数值与配置口读回一致；错误团体名不响应，错误团体名计数加 1；用户消息交付不受影响 & 37611/0\newline 926/0 & 通过 \\
SVC-02 & Set 后读回；向只读对象、错误类型、错误团体名发送设置请求；切换开关触发启动通知和认证失败通知；穿插用户消息 & 设置值读回一致；拒绝时错误字段正确，寄存器保持原值；两类通知字段正确 & 8563/0\newline 188/0 & 通过 \\
SVC-03 & 评委入口；Get、GetNext、Set、Trap 四类交换 & 四类报文字段解码正确，数值与读回一致 & \SnmpChecks/0\newline --- & 通过 \\
\end{items}

\figWaveDual[SNMP Get 请求与响应]{snmp_get_1}{snmp_get_2}{管理端 Get 请求到达代理}{代理处理后在 A、B 两网同时发出 Response}{ModelSim 仿真波形：一次完整的 SNMP Get 报文交换}{fig:wsnmp}

\section{615A 数据加卸载}\label{sec:f-615a}

\textbf{测试目的：}验证 ARINC 615A 加载与卸载流程、TFTP 块传输边界、异常重试和片内 RAM 内容，对应测试点 F4-01 至 F4-05。加载端经应用口收发 TFTP 报文；加载文件由参考程序按 615A-2 独立生成。

\begin{items}{615A 测试项}{tab:f-615a}
SVC-04 & 加载：LUI、LUR、头文件、数据文件、LUS；卸载：LND、LNR、LNS，以 \code{MEMORY.BIN} 回传。数据长度 1500、1536 字节；文件版本 A4、A3；请求端口 59、1059；穿插用户消息；注入校验错误、丢数据块、丢确认 & 见下文四项判据 & 44098/0\newline 752/0 & 通过 \\
\end{items}

判据为以下四项同时成立：
\begin{enumerate}
  \item 线路上传输的各文件经独立解码，格式正确；
  \item 状态文件的进度和最终状态正确；重试后完成，重试耗尽时正确终止；
  \item 加载完成后经配置口逐字节读出片内 RAM，与源文件一致；
  \item 卸载回传的文件与原加载数据一致。
\end{enumerate}
1536 字节恰为 3 个整块，按 TFTP 规定须再传一个空数据块结束传输；1500 字节的末块不足 512 字节，直接结束。两种长度分别覆盖这两种结束方式。

\figN[TFTP 数据块与确认]{t15_w_615a}{加载过程中的 TFTP 数据块与确认：被测端收到 DATA 块后在两网发出对应 ACK，不足 512 字节的块结束传输}{fig:w615a}

\section{RTC 时钟同步}\label{sec:f-rtc}

\textbf{测试目的：}验证 SNTP 同步功能、同步误差、时间换算和异常回应处理，对应测试点 F5-01 至 F5-06。参考检查器解码线路上每个 48 字节时间报文，核对模式字段、时间戳编码，以及响应中的 Originate 是否等于本次请求的发送时间戳。

\begin{items}{RTC 测试项}{tab:f-rtc}
SVC-05 & 空载、普通用户流量、77 条 VL 满载三种场景；空载时交换主从角色，从端时间分别超前和落后；在秒边界附近启动；停止同步后继续运行；纳秒与 NTP 小数换算 72 组向量（含零、秒末和随机值） & 报文字段与时间戳编码正确；Originate 与请求一致；秒进位正确；停止同步后时钟继续运行；换算双向一致；同步误差见表~\ref{tab:rtc-err} & 10675857/0\newline 1811/0 & 通过 \\
SVC-06 & 丢弃第一个回应；交换主从端的启动顺序 & 从端在下一周期重发请求并完成同步 & 28/0\newline --- & 通过 \\
SVC-07 & 送入 Originate 错误或过期的回应，之后送入合法回应 & 错误和过期回应不校时；合法回应只校时一次 & 34/0\newline --- & 通过 \\
\end{items}

同步误差的测量方法：每次同步完成后，经配置口同时读取两端的秒和纳秒寄存器，换算为纳秒后取差。结果见表~\ref{tab:rtc-err}。带载误差 \RtcLoad{}~ns 小于一个最长帧的占线时间 12.304~µs，满载与普通带载相同，误差不随负载增大；满载期间用户吞吐保持每方向 \Thru{}~Mbps。

\begin{table}[!htbp]\centering
\caption{RTC 同步误差（SVC-05）}\label{tab:rtc-err}
\begin{tabularx}{\linewidth}{@{}Lrr@{}}
\toprule
工况 & 同步次数 & 最大误差（ns） \\
\midrule
空载 & 20 & \RtcIdle \\
用户流量带载 & 8 & \RtcLoad \\
77 条 VL 双向满载 & 40 & \RtcLoad \\
\bottomrule
\end{tabularx}
\end{table}

\figN[SNTP 请求与响应]{t16_w_rtc}{从端发出时间请求并收到主端回应；回应中的 Originate 等于请求的发送时间戳 T1}{fig:wrtc}

\section{五功能联合测试}\label{sec:f-all}

\textbf{测试目的：}验证 5 项功能在同一对端系统、同一次运行中共用线路、应用口和配置口时的协同工作。该用例为评委入口的第一个用例，并在两种综合网表上运行（第~\ref{ch:netlist}~章）。

\begin{items}{五功能联合测试项}{tab:f-all}
SVC-08 & 依次执行：(1) 两端各发一条用户消息；(2) 向一端注入一个跳号帧；(3) 一次 SNMP Get 读取设备描述；(4) 一次时钟同步；(5) 加载 16 字节文件，读回 RAM 后卸载 & (1) 双网副本去重，各交付一次；(2) 完整性错误计数加 1，交付数不变；(3) 响应经独立 BER 解码正确；(4) 配置口确认同步完成，请求与响应时间戳对应；(5) RAM 与卸载内容均与原文件一致；结束时交付数与错误数无多余增加 & \SvcChecks/0\newline 65/0 & 通过 \\
\end{items}

独立核对的 65 项中，下载内容核对 41 项。图~\ref{fig:welinx} 为 eLinx 原生环境下运行系统级行为仿真的界面波形截图。

\figWave[eLinx 行为仿真窗口]{tb_elinx}{eLinx 原生环境运行 tb\_elinx 行为仿真界面波形截图}{fig:welinx}
